\section{Lecture 8 (September 25th)}
\begin{rmk}
Last time, we've seen how $[{\bm Q}(\alpha ):{\bm Q}]$ is a finite power of 2 for a constructible $\alpha $. Today, we'll learn about splitting fields.
\end{rmk}
\vspace{2ex}
\begin{thm}
(13.38) Let $\alpha $ be a constructible number. Then,
\begin{itemize}
\item[(i)] $\alpha \in F$, where $[F:{\bm Q}]$ is a finite power of 2
\item[(ii)] $\alpha $ is algebraic over ${\bm Q}$
\item[(iii)] $[{\bm Q}(\alpha ),{\bm Q}]$ is again a finite power of 2 
\end{itemize}
\end{thm}
\vspace{2ex}
\begin{thm}
To prove this theorem, it is enough to show that: $C=\{$numbers that have been constructed$\}\subset E$, a field, acted on by the basic operators gives $C\cup\{$newly constructed numbers$\}\subset E(u)$ where $[E(u):E]\leq 2$. For example, 
\[\{y=mx+r\}\cap \{x^2+y^2+ax+by+c=0\}\]
solves to give
\[x'=\dfrac{-(2mr+ bm +a\pm \sqrt{(2mr+bm+a)^2-4(m^2+1)(r^2+br+c)})}{2(m^2+1)}\]
By setting $u=\sqrt{(2mr+b m+a)^2-4(m^2+1)(r^2+br+c)}$, we find that both $x'$ and $y'=mx'+r$ lies in $E(u)$. It is trivial from here that the field extension has a dimension lower than or equal to 2.
\end{thm}
\vspace{2ex}
\section{Chapter 14. Normal and Separable Extensions, and Splitting Fields}
\begin{thm}
(Kronecker) Let $K$ be a field, and let $p(t)\in K[t]$ be irreducible. Then there exists an extension field that contains a zero $\alpha $ of $p(t)$. 
\end{thm}
\vspace{2ex}
\begin{proof}
Set
\[L=K[t]/(p(t))\]
and we have a field extension that contains a root of $p(t)$. Let $\alpha =t+(p(t))\in L$. Observe that 
\[p(t+(p(t)))=p(t)+(p(t))=(p(t))=0\]
and that $\alpha $ is a root of the polynomial $p(t)$ in $L$. It is always great to gain mastery in multiple notations. We can also write the above using `bar' notation as
\[\overline{p}(\alpha )=\overline{p}(\overline{t})=\overline{p(t)}=0\]

\end{proof}
\vspace{2ex}
\begin{ex}
Let $K={\bm R}$ and $p(t)=t^2+1\in {\bm R}[t]$. The field of complex numbers can be defined as
\[{\bm C}={\bm R}[t]/(t^2+1)\]
and would contain the root $\alpha =\overline{t}$. Note that
\begin{itemize}
\item[(i)] \[\overline{t}^2+\overline{1}=\overline{t^2+1}=\overline{0}\in {\bm R}[t]/(t^2+1)\]
That is, $\bar{t}$ is a zero of
\[x^2+1\in \Big({\bm R}[t]/(t^2+1)\Big)[x]\]
\item[(ii)] Note that the map $\iota:{\bm R}\hookrightarrow {\bm R}[t]/(t^2+1)$ is injective
\item[(iii)] ${\bm R}[t]/(t^2+1)$ is a field since ${\bm R}[t]$ is a PID and $t^2+1$ is irreducible in ${\bm R}$. 
\end{itemize}
\end{ex}
\vspace{2ex}
